\chapter{第三集 \quad 星斗}


\section{正文}


长江、黄河自雪域发源，宕落高原蜿蜒于群山与莽原之间，向东汇入大海，古老的中国文明在大江大河之间生长。然而北有荒漠雪山、西伯利亚荒原和蒙古高原横亘；门上有青藏高原、横断山脉阻隔；东、南皆是茫茫大海。

早期的先民基本处于一个相对独立的地貌环境，与南亚、西亚、中亚和地中海等区域相比，古中国文明与欧亚其他早期文明之间道阻且长，难以建立起直接的联系。

于是在漫长的时间岁月里，这片大地上的先祖们经风雨，历霜雪，斩斫着沿路荆棘，探索着漫漫前进之途。

终于在距今5,500多年前后，因着长期的稳定与繁荣，人烟蕃息之处越来越多，晋陕豫三省交界处这片横楔在秦岭与黄河之间的狭长平原出现了高陵杨官寨、华县泉护村、灵宝北阳平、灵宝西坡这些面积数十万甚至百万平方米的仰韶文化大型中心聚落。

初夏野蔷薇盛开在华山脚下的泉护村，绘陶青年的面前摆满了兽骨、兽牙、树枝等绘画工具以及各种可化为颜料的矿石。他将所需矿石砸碎，再放置研磨器中，用石棒研磨成粉。这是迄今发现最早的一套完整的仰韶文化绘画用具，由陶杯、石盖、石砚、研磨棒和氧化铁颜料组成。石砚臼窝内壁及砚面上还残留红色颜料的痕迹。

在青年工匠的笔下行云流水地绽开一朵朵花，经过火的淬炼，黄底黑彩的陶器便能再现5,500年前那个夏日人们所见过的绚丽，闻过的香气，陶上怒放的野蔷薇将不再步入花开花谢的轮回，而将永世盛开。这类绘有花瓣纹样的陶器，被今天的考古学家命名为仰韶文化庙底沟期彩陶。

仰韶文化历时近2千年，庙底沟期是它的中期，因河南三门峡庙底沟遗址而得名。花卉纹是庙底沟彩陶最具代表性的纹样，除此之外还有回旋勾连纹、豆荚纹等。

在低饱和度的黄色或红色素地上，以黑彩、白彩或红彩描绘纹饰，通过连续、反复、对称、共用的手法，实现了写实与抽象的意向表达，阴纹与阳纹的相互映衬，平视与俯瞰的视觉效果。

曾经的仰韶文化早期各地虽然都流行彩陶，却各具区域特色，到庙底沟时期从关中平原山西南部到河南西部，彩陶的纹饰模式高度近似，这跨越距离的呼应却是因何而起？

考古学家发现史前人群的分野往往不仅以血统为要素，更以生活方式或文化传统当作彼此辨识与认同的密码。

庙底沟人原以旱作农业为主，也逐渐接纳了稻作农业的升级方式，这就是不同人群接触交流甚至融合的重要迹象。

不仅如此，籍由成熟早作农业人群的扩散和流动，庙底沟彩陶的影响先是到达了中原河洛地区，再以整个黄河中上游为起点，挟风行之势，向着北抵阴山，东到大海，西达甘青，南至长江的更广阔天地辐射。

向东山东海岱地区大汶口彩陶的叶片纹、花瓣纹和旋纹是自庙底沟彩陶掀来的风潮；向西传承陕西关中的甘肃彩陶继续朝着更远的陇西和青海东界传递影响；在稍晚的马家窑文化时期，甘青地区将发展成为新的彩陶文化中心；北溯汾河和南流黄河两岸的谷地、高原而上，庙底沟彩陶为内蒙古中南部及东北的红山文化彩陶注入了新的生命力；南延汉江而下，庙底沟彩陶在长江中游的豫南、鄂西北引发追捧后，继而流播至湖南、浙江等地。

古老的华夏大地由此拉开第一次大范围交流、大规模互动的序幕。在没有广告效应和发达物流的史前时代，器物传播依靠的是人与人之间的传递。

在各地都已发展出成熟陶器传统的那个时代，由庙底沟掀起的彩陶艺术浪潮能够波及这样广阔的区域，应该代表着远距离交流的频繁以及以彩陶为基础的文化共识，这或许就是文化意义上的最早中国的雏形。

在后世的《史记》中保留了有关黄帝这位华夏共祖的传说和记忆。据说他的足迹东至大海，西至陇右，南至长江，北至幽燕，与庙底沟彩陶的影响范围惊人的重合。

在古汉语中"花""华"同音，"华"的本意为"花"。金文中的"华"字就是花朵加上花蒂的样子，因此考古学家苏秉琦将庙底沟彩陶中常见的纹样喻为"华山玫瑰"，而把创造庙底沟之花的人群称为最初的"华人"。美丽的花卉纹似乎可以解读为华夏之"华"的最初由来。

土地肥沃河流纵横的关中平原是仰韶文化绽放的核心区域之一。泾河与渭河交汇处的杨官寨遗址面积达百万平方米，考古发掘揭露出由大型环壕、西门址、中心水利系统、东区墓地等组成的规划有序的村落布局。东区成人墓地已经发掘500多座墓葬，显示了聚落繁荣发展的面貌。

杨官寨遗址出土的特殊器物，如镂空人面漆座、朱砂人面饰陶器、动物纹彩陶盆等，可能与聚落进行的重大祭祀活动有关。表明当时已经拥有相对复杂的政治组织，为社会向更高一级发展奠定了基础。

秦岭余脉绵延向东直指中原，黄河南下穿过龙门转折东流，河南西部灵宝的铸鼎原，南依秦岭，北面大河，是仰韶文化的另一个核心区域。其中西坡遗址是经发掘确认面积最大的聚落之一，它的中心是4,000多平方米的广场，广场四角各有一处大型建筑。

室内的火塘中烈焰正燃，一场庄严的成年礼即将在这里举行。一位青年站在大房子之外难以激动的心情，只等完成这场隆重的仪式后他就能成为真正参与族中事务的一员，这是他自孩提时代起就抱有的梦想。

大房子的建造曾是村中最大的工程。不包括砍伐、加工和运输树木等工序，单单建造房子本身就让百余名壮劳力施工足足3个月，彼时还是稚龄的他充满好奇总想一探究竟。建造这样的大型建筑，先民采用了相当复杂的工艺。先要开挖房基坑并垫土夯打，在基坑周围挖出墙基槽，槽中竖立粗大的木柱，用夯土填满基槽使木柱牢固。通过夯打与夯填形成内外墙体，此后设置室内柱础石，挖门道、火塘，栽立室内柱，再用不同种类的黏土和夯土铺垫坑底，以石灰膏泥抹出平整的地面和墙壁，最后依托木柱搭建起回廊、门道以及屋顶。这类大型建筑的出现体现了聚落上层日益强大的组织能力，是庙底沟社会发展到新阶段的又一有力证据。

当年的西坡少年曾趁着众人休息时偷偷溜进大房子里，眼前所见令他无比吃惊，就像有几回夕阳余晖染红天地河的景象，墙壁和眼前的柱子都被刷成了红色。恍然之间他发现位于中央的一根柱子还没上完色，他知道这是大房子里最神圣的区域，他忍不住蘸了颜料就要去刷柱子，却被门外突然的响动吓个正着。

未成年的族人是不允许进入这样的重要场所的，可出乎意料的是面对这擅自闯入的少年，族长却没有一丝恼怒，他轻轻扶起眼前的孩子，示意他应该从上往下刷。少年迟疑地探出手臂却一不小心将红色的颜料甩落在族长的脸上。

当这回忆沉淀，步伐就愈加坚定，那个当年偷溜进来的孩童已经长大成人。穿过狭长的门道绕过烈焰腾腾的火塘，进入族人们簇拥着的庄严空间。这铺天盖地、动人心魄的红，红色的立柱、红色的墙壁、红色的地面仍一如当年．..

这座大房子曾见证过无数重要的时刻，西坡和附近小聚落的居民们在特定的时刻总会汇集在此，举行各种各样重要的仪式。神圣的宗教典礼、社群的公共仪式、盛大的婚礼、集体的宴饮、狂欢的节庆。今天的主角是这位西坡青年，在成年礼之后他就要为族中承担起更大的责任。

后来考古学家在西坡遗址的中心发掘出这座大型半地穴式建筑遗址。房址的室内面积为204平方米，加上外部回廊，总占地面积超过500平方米，是迄今发现的那个时代最大的单体建筑。

庙底沟时期许多仰韶文化聚落中的大房子规模宏伟、工艺复杂。他们往往位于聚落中心，宽敞明亮，结构考究，地面层层加工铺设，墙面多施彩绘，想必不是普通的住所而是整个聚落的公共聚会场所。

当屏息静气，仿佛已坍塌的房屋，逆转时光再度矗立，还能听见某个首领或慷慨或舒缓的陈辞，还能听见聚在这里的人们，那热闹的喧哗。

随着彩陶的广泛传播，仰韶文化的社会内部和外部均建立起空前紧密的交流网络。中原和位于东部的海岱地区便有着频繁而持续的互动。

父辈们曾经赶赴大海之滨的壮阔旅程，远方来客不断诉说的奇风异俗，频频拨弄了青年的心，踌躇满志的西坡青年决定在成年礼的第二天就向东出发。

西坡青年沿着黄河一路向东，从西坡到大汶口的距离是600多公里，如果按照每小时行进5公里，每天行走6小时计算大约需要20多天。青年朝行夜宿，路上看过几十轮日升月落，穿过辽阔的平原行至泰山脚下，终于到达了大汶口聚落。

当黄河中游经仰韶文化的整合而掀起风起云涌的变革时，下游也正由大汶口文化引领着书写下气势恢宏的新章。大汶口文化自岱宗泰山和渤海黄海之间的古济水与泗水、汶水流域崛起，一统今山东省全境及皖北、苏北、豫东等地大有风雷之势。

西坡青年受到了大汶口人的热情款待，参加了令他生平难忘的一场盛宴。宾主面前陈列着各种食器和酒器，陶豆中装满了肉干和水果，陶鼎内盛满了金黄的黍糜，精致的高柄杯中斟满了清香的美酒．..

大汶口社会上层豪气而富有，热衷于举办盛宴款待宾客，这也是他们展示财富和威望，广泛交友，扩大往来的媒介。

"好客"或是远古的基因传延至今。宴饮的奢华并未让西坡青年过分艳羡，席上尊卑有序的饮酒进食之礼，却让他格外留意。

考古研究者发现大汶口文化的大型墓葬往往随葬大量食品以及与饮食相关的陶器。大汶口遗址2005号墓的棺椁之间有数组陶器，一组是饮器，如高柄杯、觚形杯；一组是食器，如鼎、豆、壶、三足钵、三足碗等。椁外放置了成套食器，器内多盛有肉食。

大汶口文化不同区域的大型墓葬，几乎都采用了类似葬仪。这说明即便地跨山东、苏北和皖北，各聚落的社会上层之间仍保持了频繁的交流。

相似的考古发现背后是一致的文化认同和社会秩序。以饮食宴享和成套器具为媒介构建的礼仪，则是大汶口文化权利与地位的象征。席间那些绘有花瓣的彩陶瓶，让西坡青年觉得格外亲切，这他熟悉的故乡风物，是他的先辈与这东方之族曾远程往来的凭据，也是仰韶彩陶之风被初兴的大汶口文化效慕接纳的明证。

主人一杯杯地劝，客人便一杯杯地饮。交谈时西坡青年注意到，在座不少大汶口人的上颚都缺了一对侧门齿，这种拔牙之风乃是族中世代相传的习俗，或许是要用对于痛苦的忍耐来显示勇气，宣告成年。

那天的盛宴上除了西坡青年，还有另一位来自异乡的年轻人。在一片欢声笑语中，那位异乡人却显得拘谨而沉默，他是从淮河流域北上的凌家滩人。

大汶口人夸耀着丰盛的食物和财富，强调着世俗社会的秩序和礼仪，而他所来自的凌家滩社会却更注重宗教精神的引领。

现实的欢乐和规仪并不能打动凌家滩青年，唯一让他心有戚戚的是主人手边的一件陶豆。红陶之上以白彩描画的八角星纹格外显眼。在他的家乡八角星纹是神圣的标志，这难道是一种巧合吗？

凌家滩在安徽巢湖一带，与同在长江下游的崧泽人、北阴阳营人一样都有着尚玉之风。大部分玉器都是非实用性的装饰品或仪式用品，这意味着那个举全族之力只为生存的艰险岁月已经过去了。人们产生了除温饱之外的更多需求，社会已能释放出部分人力来专门制造这些"无用之用"的器物。

那时的玉器制作工艺已经相当成熟，开坯时广泛采用线切割或片切割的技术；开孔时则采用了管钻技术，坚致的玉能被人们自如地雕琢成想要的样子。

凌家滩青年的父亲是部落里最好的玉匠。父亲的那双手简直有着无穷魔力，他在童年时总是如痴如醉地观看父亲切割、穿凿、琢磨，一块块看似不起眼的石料经他之手便能化为精巧的美玉。乍见这片玉板之上的"八角星"，他蓦地怔住。回过神后不断地追问父亲其中的含义，父亲只是笑而不语，让他长大后自己去寻找答案。

父亲曾为族里离世的祭司和首领制作过数件玉人。玉人或坐或立，双臂紧贴身体，弯曲上举于胸前，臂上穿戴着成排的镯饰，似是在虔诚祈祷。这种姿势颇含仪式意味，是族人们在特定时刻的程式化肢体语言，也是祖先入葬时的造型。

最令青年赞叹的还是父亲制作的一件玉鹰，那鹰展翅飞翔，仰首侧视，展开的双翼各如一只猪首，敞露出胸腹部赫然的八角星纹，说不出的神秘莫测。

凌家滩人在玉器中寄寓着他们的信仰观念，在日益分化的社会里也以玉器作为展现财富和权力的重要载体。凌家滩墓地建在大型祭坛之上，逝者按照身份的区别被葬入不同的区域，拥有多寡不等和种类不一的随葬品。

07M23号墓位于墓地南侧的中央区域，墓主人的遗骨因数千年环境腐蚀已经荡然无存，但随葬的300件玉石器却印证了他的首领身份。

考古学家通过详细的现场记录和反复摆放分析，复原了当时复杂又严格的礼俗。葬具的最底层用穿孔石钺横向摆成数道"横木"的样子，其上就是七排严丝合缝整齐排列的石、石凿，刃部朝向墓主人的脚步，密密麻麻地砌出一座"平台"。在长2.65米，宽0.8米的葬具内，墓主人的遗体就曾安放于这象征尊崇身份的基底之上。

一种特别的"花石头"被专门用来制作大孔弧形刃部的特殊石钺，在胸口部分放置两排六件，然后沿着墓主身形的中线，整齐往脚步首尾相连排成一线，这些石钺显然没有使用过，入葬时也未装柄，而是摆成一个"T"字形，成为墓主置身的第二层基础。

玉器是身上最重要的配饰，墓主脖梗上挂着一套穿缀繁复的玉璜组佩，两臂上分别穿戴10件玉镯，腰腹部则穿挂一套内置玉签的特殊礼器，造型如龟壳互扣。穿戴整齐的墓主人身上还要继续用石钺和玉环、玉玦来铺陈。

石钺从材质和技术上都能看出三件一组摆放的规律，分别压在身体两侧，尺寸从大到小。成套制作的玉环、玉玦密密地挂坠在葬具的两端，墓坑内散落的玉环可能是最后才撒在葬具上的。除了随身佩戴的玉饰外，墓中出现的大量的玉石器都是为了这样的仪式活动而定制的。用特殊的玉石资源，结合专门技术生产的手工业"奢侈品"由此成为了早期中国体现"身份"的重要内容。

在敷土填满整个墓坑的最后，一件重达88千克的玉石猪被特意摆放在墓土中。这件长 72厘米，宽32厘米的玉石雕是目前中国史前发现尺寸最大的作品。整个葬礼奢华又严格，是凌家滩聚落体现个人地位的重要仪式。

显贵者们被埋葬在墓地的南区，青年的父亲去世后跟其他族人一起葬在西北的墓区。编号98M20的墓葬中除了凌家滩常见的玉石器外，还出土了111件加工玉器后所剩下的玉芯和材料，正是这些特殊的随葬品为考古学家推测墓主人的身份和职业提供了线索，他可能是一位专职的玉工。

一个好的匠人不仅要有精湛的技艺，还需有内心的充盈。父亲的这句话他不曾忘怀。多年前父亲曾听人说起遥远的北方有一处圣地，那里的人们也善于琢玉，并通过仪式中复杂精美的玉器来表达信仰与寄托。父亲心向往之，但终未成行。如今对父亲最好的纪念就是走他想走的路，寻找那跨越生死与时间的终极意义。

是时候离开华服盛宴的大汶口了，凌家滩青年与西坡青年道别，又向好客的主人家辞行，坚定地向着圣地红山行进。

凌家滩青年翻过泰山，沿着茫茫的海河平原向东北前行，当他穿过辽西走廊来到巍峨的大青山脚下便进入了红山文化的核心范围。努鲁儿虎山深处，牛儿河畔的山谷间有一道漫长的山梁，当地人依河名，称其为牛河梁。穿过起伏的山峦，圆坛、神庙和一座座积石高家，在林端时隐时现。

红山人对精神文化的需求高于一切。公共权力自宗教层面产生形成了层级严明的体系，而牛河梁拥有着其中规格最高的、完整的坛、庙、冢祭祀体系。于山梁之上修建山台，西南建筑群采用中轴对称布局，轴线南端直指远处的猪首山山峰，如此规划可谓苦心孤诣。然而红山人的空间视野远不止于此，在牛河梁周围近50平方公里的范围内，他们充分利用山形地势的变化，勾勒出宏大的布局。

云雾深处，红山女巫的身影若隐若现。凌家滩青年亦步亦趋，登上主梁。沿着漫长的山脊穿越掩映的密林，青年来至女神庙前。

这是北方史前时期流行的半地穴式建筑，但采用了较少见的多室结构，平面呈"亞"字形，总面积约 75平方米。

走进狭长的门道神秘庄严的气氛扑面而来，神庙之内装修极为考究，地面经火烘烤，木骨泥墙上装饰有仿木建筑构件，墙面还绘有赭红、黄白相间的彩色几何纹样。相当于真人三倍、两倍和等同的数尊女神塑像，伫立于北主室内，女神头像眼嵌玉石、齿贴蛙饰。

在那个时代产妇和婴儿的死亡率很高，人们哀劬母亲诞育子女的不易，也期望人口能顺利繁衍，将之投射到信仰中便是对于女神或母神的崇敬。

走出女神庙时已是黄昏，落日映红牛河梁的峦。山间散步的积石冢是族中先人们的安葬之处，也是后代族人崇奉先祖的场所。

一冢之中往往有多座墓葬，因墓主身份的不同而被分别安置在大小各异的石砌葬具之中。在牛河梁的积石冢内，玉器几乎是唯一的随葬品，但数量并不算多。大型墓葬中一般也不超过10件，组合也不十分固定，与材质单一数量不多的随葬品形成鲜明对比的却是积石冢那异常巨大的规模。

这种不相称恰恰反映出"唯玉而葬"现象的内涵，那便是玉器无关财富的占有，也不刻意显示等级的差别，而是与特定的象征功能相关，或许是作为通神的工具。

5号地点1号家的这座中心大墓，墓主头部两侧对称摆放着玉璧，腹部放置勾云形器，双手各握一件玉龟。

2号地点1号家的南侧墓群中，埋葬着一位富有的成年男性，随葬玉器20件，斜口筒形器、兽面牌饰、勾云形器、玉龟等等。这些明显具有观念意识的器物折射出红山文化积石冢和葬礼仪式中浓厚的原始宗教特色。

双手置于胸前的玉人、蜷曲的玉龙、斜口筒形器、玉龟……不少红山文化玉器与凌家滩文化的玉器颇具相似之处。环绕积石冢的无底彩陶筒形器，器表施有勾连涡纹，又让人联想到庙底沟彩陶的影响。这些发现似乎证实着5,000多年前黄河中下游、长江中下游和辽河流域等地的社会上层之间已经存在远距离的物质和精神文化交流。

青年望向不远处的圆形祭坛，红山女巫已经站立在祭坛中央。祭坛由三个同心圆石圈构成，每一圈的边缘都竖立着多棱形的花岗岩块，最内圈的整个地面铺满白色石灰石。由外至内的三圈或许分别象征太阳在冬至、春分、秋分以及夏至运行的轨道。这样的结构特征与唐长安城南郊圆丘、明清天坛较为相似，体现出古人对天地方圆早有了超越时间的共通之理。红山女巫默诵祷词，双手置于胸前，虔诚而肃穆地展开双臂，如同展翅的雄鹰将飞升而上。凌家滩青年目睹这庄严而神圣的情景，内心如有和鸣，豁然洞明，唯有将心灵与精神托付于那伟大的自然规律，敬畏天地，方能得到持久的安宁。

西坡青年在大汶口停留了一段时间后踏上了回家的路。回顾这段旅程，大汶口的礼仪和物质生活令他印象深刻，而那些与他所在的中原既有联系也有区别的种种是一段更加难以磨灭的记忆。

他迫切想把此行的所见所闻告诉部落里的每个人，见识了更大的世界是一个青年成人的骄傲。当青年滔滔不绝地讲述起这段旅程的种种，是否会引起族人们长久的惊叹？考古论证止步于此，难以描述生动的细节，但这位经过漫长旅行洗礼的年轻人或许终将成为西坡村落新一代强有力的领导者。

生活似水，又流过了几代人的光阴。西坡村落也在悄然发生着变化。自先辈肇始的对外交流，逐渐扩容了本地文化的基因。

考古学家从这一时期的大型墓葬中辨识出了仰韶先民开创独特社会发展道路的足迹。5,000多年前的盛夏，在黄土塬上的西坡聚落中正举行着一场盛大的葬礼。这是显示社会身份，维系社会组织和秩序的重要仪式。

精细的考古发掘可以让我们推想那场葬礼的部分细节。整个葬礼大致是从遗体处理开始的，首先要清洗逝者的身体和头发。西坡墓地中有些墓主头戴发簪，应该就是清洗头发以后插上的，待清洗过程完成后，人们用织物小心地包裹逝者的身体。在下葬的日子，长长的送葬队伍从中心广场出发，缓缓南行，由聚落的南门跨过壕沟来到地势高亢的墓地，应该会有幽咽的陶损吹响吧。

苍翠的秦岭近在眼前，回首北望黄河隐没在远山之间，在众人肃穆的注视下，包裹着麻布的墓主人被轻轻放置在幽深的墓中。

人们仔细地摆设随葬品，然后用木板封盖墓室和脚坑。

经过多学科研究，我们对这位身高约 1.65米年龄在35岁左右的男性墓主有了更深的认识。他的右侧肋骨中部有骨折错位愈合现象，生前可能因参加狩猎、战斗或竞技而受伤。寄生虫专家提取他盆骨内的土样进行分析后，发现了大量因食用猪肉滋生的寄生虫卵，而猪肉是那个时代奢侈的美食。

我们或许永远无法知道这位男子姓甚名谁，但有一点可以明确，他是庙底沟社会等级化的代表，是距今5,000多年前贵族阶层中的重要一员。

值得注意的是他下颌的两颗门齿在生前就被有意拔除，这是大汶口人流行的风俗。他的墓中还随葬着东方风格的大口缸或许随着大汶口文化的发展壮大，连带着当地的一些饮食器具已逆着彩陶东去的路途，反向输回到中原地区。

庙底沟的先民建造了大型聚落、大型公共建筑以及大型的墓葬，但考古发现中既无奢华的随葬品，也无浓厚的宗教气氛，而用来凝聚族群和巩固权力的大墓葬仪却可以看到跨区域远距离文化交流的借鉴和影响。

简朴实用与多元开放成为中原地区社会文化的特质，对于此后中国文明的形成与发展有着深远影响。

距今6,000年以来，随着庙底沟文化彩陶的扩散，大汶口文化的西渐以及东方各玉文化中心的形成，黄河流域、长江流域和西辽河流域间的文化互动交融日趋紧密。

这满天星斗竞相闪耀的大地，塑造了中国的地理核心。从此区域文化各有所长，发展并进，积累起中华文明多元化的深厚土壤。

区域间日趋深化的交流与互动为一体化的将来积蓄着力量。那一群群阔步山川间的年轻人正如我们的青春中国，已昂首伫立在文明盛放的前夜。
